\scanpage{136}\mspage{69}

de ce dernier, on trouve une décomposition
\[
F=\coprod_{i\in\Pi}F_i,\qquad A=\coprod_{i\in\Pi}A_i\tag*{(\deleted{$\unsure{T}$}24)}
\]
les $F_i,A_i$ non vides et mutuellement disjoints.

On exige alors (cf p. 53)

(25)\qquad (T 3)\qquad $\forall i\in\Pi$, $(S,F_i)$ est une structure \uline{triangulée} \uline{stricte}\inserted{ $P_i$}, i.e. $\bigcup_{f\in F_i}f=S$, et on a connexité du contour local en $s$ défini par $P_i$,

cf T3 a) b) p. $53'$ pour une explicitation.

Les autres axiomes peuvent s'exprimer, en disant que la famille $(F_i)$ des structures triangulées sur $S$ forme exactement une classe d'équivalence de structures triangulées symplectiques précongruentielles.

% slanted
\margin{En \unsure{travaillant|traduisant|dimension} avec les structures triangulées $(S,F)$ de $=$, on trouve que (25 bis) $\left\{\begin{aligned}&G\text{ est simplement transitif}\\&\text{sur l'ens. des ``\uline{repères}''}\\&\text{de }(S,F)\end{aligned}\right.$ donc il est transitif sur $\vec F,F,\vec A,A,S$, le stabilisateur d'un él. $\vec f$ resp $f$ resp. $\vec a$ resp $a$ resp $s$ est isomorphe à $\mathbb Z/3\mathbb Z$ (can.), $\unsure{\mathbb D|D}_3$, $\mathbb Z/2\mathbb Z$, $\mathbb D_{\unsure{2|{}}}^{2}$, et est \unsure{isom|isom.} (pour $s$) à un produit $\frac12$ direct}

Il faut en arriver enfin, progressivement, à des conditions plus strictes encore, pour en arriver enfin aux \inserted{structures} \deleted{qu'on aura envie d'appeler} ``congruentielles'', isomorphes à celles déduites des modules libres de rang 2 sur des $\mathbb Z/n\mathbb Z$. Pour l'instant, l'attention principale \deleted{\unsure{est}} sera sur les phénomènes liés au cas d'un $\Pi$ non réduit à un point, i.e. d'un $Z$ qui n'est pas réduit à $\{1\}$ (ou à $\{1,\unsure{\underline{\circ}|\underline{-}|\varepsilon}\}$, dans le cas

% archive sheet number: 135
